% =============================================================
% §5 实验评估
% =============================================================
\section{实验评估}
\label{sec:eval}

\subsection{实验设置}
\label{sec:eval-setup}

\textbf{仿真平台。}所有实验均使用 \S\ref{sec:simulator} 的仿真器，其以里程碑 M1--M4 实现五项 AIGC 物理。

\textbf{硬件配置。}1 台云级服务器（200~TFLOPS，128~GB 显存，等效 A100-80GB）+ 1--7 台异构边缘节点（10--50~TFLOPS，16--64~GB 显存，对照 Jetson AGX、T4、边缘 A100 部署建模）。云--边链路网络延迟 30--65~ms，边--边链路 8--23~ms。

\textbf{工作负载。}采用与 Azure LLM Inference Trace 2023~\cite{azure2023llmtrace} 公开统计对齐的对数正态请求长度分布：prompt 长 $\sim \text{LogNormal}(\mu{=}5.5, \sigma{=}1.0)$，输出长 $\sim \text{LogNormal}(\mu{=}4.5, \sigma{=}1.2)$。请求到达为速率 $\lambda$ 的 Poisson 过程。每个请求在三种模型（LLaMA-7B、LLaMA-13B、LLaMA-70B-INT8）中均匀采样，模型参数标定至 vLLM 与 TensorRT-LLM 的基准数据。

\textbf{基线（11 种调度器）。}对比覆盖广泛：
\begin{itemize}
\item \textbf{RoundRobin（RR）}---无负载感知基线。
\item \textbf{LeastLoaded（LL）}---负载感知贪心。
\item \textbf{ShortestQueue（SQ）}---队列感知贪心。
\item \textbf{HEFT}~\cite{topcuoglu2002heft}---标准 DAG 调度基线。
\item \textbf{GA}~\cite{holland1975adaptation}、\textbf{PSO}~\cite{kennedy1995pso}---元启发式搜索。
\item \textbf{NSGA-II}~\cite{yu2025llmrouting}---多目标遗传搜索（非支配排序 + 拥挤距离），在逐窗口分配上进行，两个目标与本文两轴对应（完成时间代理、能耗估计）；执行最终 Pareto 前沿的 knee point 解。
\item \textbf{A3C-R2N2}~\cite{tuli2022a3cr2n}---RL 基线，采用异步优势 actor-critic 与 GRU+残差编码器，但使用\textbf{通用（非 AIGC）状态与奖励}。
\item \textbf{GNN}---本文调度器的 graph-neural-network 变体，将 AIGC 物理编码为任务--服务器二部图中的\textbf{边特征}，与 Decima~\cite{mao2019decima} 类似。
\item \textbf{CS-UCB}~\cite{yang2024perllm}---PerLLM 风格的约束满足 UCB bandit，最接近的先行工作（\S\ref{sec:related}）：在可行 arm 中按最大 UCB 逐请求选择服务器，奖励为 $-E_{\text{est}} + \lambda f(y)$（负能耗估计加归一化约束裕量）。
\item \textbf{RL（本文）}---本文的 AIGC 感知 PPO 调度器，含显式 AIGC 状态特征与奖励整形。
\end{itemize}

全部 11 种调度器均以 Python 原生实现，对接仿真器统一的 \texttt{scheduler.schedule()} 接口；未复用任何外部调度器库。NSGA-II 基线复现~\cite{yu2025llmrouting} 的多目标云边路由设定。HEFT 遵循原始的 upward-rank 公式~\cite{topcuoglu2002heft}；A3C-R2N2 复现~\cite{tuli2022a3cr2n} 的 GRU + 残差跳连网络；GNN 变体采用 Decima 风格的图编码~\cite{mao2019decima}；CS-UCB 在本文的逐任务粒度上复现 PerLLM 的约束满足 UCB~\cite{yang2024perllm}，以 (model, phase, server) 逐组维护的 arm 统计对应其按服务的个性化。表~\ref{tab:hyperparams} 列出各调度器的超参数与搜索/训练预算；所有实现随本文开源发布。

\begin{table}[t]
  \centering
  \caption{调度器实现与超参数。RR/LL/SQ = RoundRobin、LeastLoaded、ShortestQueue。}
  \label{tab:hyperparams}
  \small
  \setlength{\tabcolsep}{4pt}
  \begin{tabular}{@{}l@{\hspace{6pt}}p{0.64\columnwidth}@{}}
    \toprule
    Scheduler & Configuration \\
    \midrule
    RR / LL / SQ & 无参数贪心规则 \\
    HEFT & upward rank + earliest finish time；无参数 \\
    GA & 每调度窗口 population 30、50 代；交叉 0.8、变异 0.1、精英保留 0.1 \\
    PSO & 每窗口 30 粒子、30 次迭代；惯性 $w{=}0.7$、$c_1{=}c_2{=}1.5$ \\
    NSGA-II & 每窗口 population 30、30 代；交叉 0.8、变异 0.1；执行最终非支配前沿的 knee point 解 \\
    A3C-R2N2 & GRU 单元 + 残差跳连（隐层 128）；Adam $3{\times}10^{-4}$、$\gamma{=}0.95$、熵系数 0.05、value 系数 0.5；每 32 个 transition 更新一次（单 epoch、无裁剪，按~\cite{tuli2022a3cr2n}） \\
    GNN & 任务--服务器二部图上的 1-hop 注意力消息传递（隐层 64）；与本文相同的 PPO 配方与预算 \\
    CS-UCB & arm = (model, phase, server)；以准入检查作可行性过滤；探索系数 $\delta{=}0.5$、约束权重 $\lambda{=}1.0$（在 $\delta \in \{0.1,\dots,2\}$、$\lambda \in \{0.5,1\}$ 上网格调参） \\
    RL（本文） & MLP $256{\to}128{\to}64$；PPO 裁剪 $\epsilon{=}0.2$、GAE $\gamma{=}0.95$、$\lambda{=}0.90$；Adam $3{\times}10^{-4}$；熵系数在预训练后 $0.05{\to}0.01$；每 32 个 transition $\times$ 4 epoch 更新 \\
    \addlinespace
    \multicolumn{2}{@{}p{0.97\columnwidth}@{}}{\emph{训练预算（所有学习型调度器）：}20 个预训练 episode + 30 次随机 warm-up 派发；评估期间在线更新持续进行。} \\
    \bottomrule
  \end{tabular}
\end{table}

\textbf{指标。}沿用标准 LLM 服务基准：\slo{} 达成率（满足 TTFT $\le$ 2.0~s 且 TPOT $\le$ 100~ms/token 的比例）；TTFT 与 TPOT 百分位；goodput；总能耗；以及每 token 能耗（J/tok），采用 \S\ref{sec:simulator} 标定的功率模型。

\textbf{统计方法。}所有实验使用 30 个独立随机种子；报告均值~$\pm$~标准差，并以双侧 Mann-Whitney U 检验显著性（仿真输出非高斯）。阈值：$*$ $p<0.05$，$**$ $p<0.01$，$***$ $p<0.001$。

\subsection{主结果：Pareto 占优}
\label{sec:main-results}

我们在典型配置上评估全部 11 种调度器：5 台服务器（1 云 + 4 异构边缘）、Poisson 到达 $\lambda{=}2$ 请求/秒、100 个任务单位（约 50 个推理请求）。结果按 30 次运行取均值。

\textbf{核心结论。}本文 AIGC 感知 RL 调度器同时取得最高 \slo{} 达成率\textbf{与}最低每 token 能耗，在 (\slo, 能效) 平面上\textbf{按均值表现 Pareto 占优全部 10 个基线}（表~\ref{tab:main}，图~\ref{fig:pareto}）。

\begin{table}[t]
  \centering
  \caption{典型配置（edge${=}5$、$\lambda{=}2$ 请求/秒、$N{=}30$）下的主对比。Mksp.\ = makespan（s）；E = 总能耗（kJ）；\etok{} = 每 token 能耗（J）；TTFT P95 单位为秒。粗体 = 列最优。}
  \label{tab:main}
  \scriptsize
  \setlength{\tabcolsep}{3pt}
  \begin{tabular}{lrrrrr}
    \toprule
    Scheduler & Mksp.(s) & \slo$\uparrow$ & E(kJ)$\downarrow$
              & \etok(J)$\downarrow$ & TTFT P95 \\
    \midrule
    RoundRobin    & 306.2 & 0.163 & 64.0 & 7.29 & 174.6 \\
    LeastLoaded   & 125.5 & 0.239 & 26.8 & 2.94 & 30.8 \\
    ShortestQueue & 123.8 & 0.234 & 26.9 & 2.96 & 25.4 \\
    HEFT          & 126.7 & 0.169 & 25.9 & 2.83 & 35.7 \\
    GA            & 132.4 & 0.225 & 26.9 & 2.94 & 30.5 \\
    PSO           & 129.6 & 0.235 & 25.6 & 2.81 & 41.0 \\
    A3C-R2N2      & \textbf{120.6} & 0.214 & 26.7 & 2.93 & 28.8 \\
    GNN           & 125.5 & 0.210 & 26.7 & 2.93 & 32.2 \\
    CS-UCB        & 129.9 & 0.193 & 26.0 & 2.84 & 37.3 \\
    NSGA-II       & 126.2 & 0.289 & 25.5 & 2.79 & 32.0 \\
    \textbf{RL（本文）} & 122.2 & \textbf{0.302}
                       & \textbf{23.7} & \textbf{2.61}
                       & \textbf{25.2} \\
    \bottomrule
  \end{tabular}
\end{table}

本文 RL 调度器在除 makespan 外的每项指标上均为最优；makespan 一项由 A3C-R2N2 略早完成（120.6~s vs.\ 122.2~s，差距 1.3\%）。makespan 不在本文优化目标之列（\S\ref{sec:sched-problem}），该差距也不影响 (\slo, \etok) 平面上的 Pareto 关系。

\begin{figure}[t]
  \centering
  \includegraphics[width=\columnwidth]{fig4_pareto.pdf}
  \caption{典型配置（edge${=}5$、$\lambda{=}2$ 请求/秒、$N{=}30$）下 \slo{} 达成率 vs.\ 每 token 能耗的 Pareto 图。“理想”角位于左上。本文 RL 调度器单独占据该区域。单目标基线聚集在右下且 \slo{} 明显更低；最接近的挑战者 NSGA-II 在两个轴上均落后。RoundRobin 是图外的极端离群点。}
  \label{fig:pareto}
\end{figure}

\textbf{显著性检验。}表~\ref{tab:significance} 报告本文 RL 调度器对 6 个最强基线的 Mann-Whitney U 检验。

\begin{table}[t]
  \centering
  \caption{对 6 个最强基线的 $p$ 值（双侧 Mann-Whitney U，$N{=}30$）。粗体 = $p<0.05$。}
  \label{tab:significance}
  \scriptsize
  \setlength{\tabcolsep}{2.5pt}
  \begin{tabular}{lcccccc}
    \toprule
    Metric & PSO & SQ & A3C-R2N2 & GNN & CS-UCB & NSGA-II \\
    \midrule
    \slo{} 达成      & \textbf{0.045}$^*$ & \textbf{0.046}$^*$
                     & \textbf{0.006}$^{**}$ & \textbf{0.006}$^{**}$
                     & \textbf{0.002}$^{**}$ & 0.64 \\
    总能耗           & 0.36 & 0.07 & 0.09 & 0.10 & 0.21 & 0.27 \\
    每 token 能耗    & 0.11 & \textbf{0.010}$^{**}$
                     & \textbf{0.028}$^*$ & \textbf{0.016}$^*$
                     & 0.066 & 0.12 \\
    \bottomrule
  \end{tabular}
\end{table}

\slo{} 达成率维度上，本文调度器对全部 5 个单目标基线均显著更高（$p<0.05$，其中对两个 DRL 基线与 bandit 均为 $p<0.01$）。能效维度对 6 个中的 3 个显著更优（PSO，$p=0.11$，与 CS-UCB，$p=0.066$，方向一致但未达 $p<0.05$）。在该工作点上唯一未能与本文调度器统计分离的基线是 NSGA-II（\slo{} $p=0.64$、\etok{} $p=0.12$）：本文在两个均值上均更好，但在 $N{=}30$ 下差距处于噪声之内；下文将回到这一对比---包括其甜区显著性与代价不对称性。总能耗维度展示出一致的 $\le 10\%$ 优势但未达单独显著，反映了 makespan 与能耗之间预期的相关性。

\textbf{为何 AIGC 感知 RL 能耗更低？}反直觉的是，本文调度器---尽管可以调用高功率云端资源---比默认偏向边缘的基线消耗\emph{更少}的能量。我们将其归因于三种学到的行为：（i）\emph{更短 makespan 放大空闲节能}（集群空闲功率约 150~W）；（ii）\emph{冷加载规避}---将同模型请求集中放置，消除了让 GPU 通电却不产 token 的加载时段；（iii）\emph{兄弟亲和}---decode 与其 prefill 同址放置，避免了“通电空转”的 KV 传输区间。\textbf{当能耗作为协同目标时，这些行为使 AIGC 感知设计呈现 Pareto 占优。}

\textbf{与 A3C-R2N2 对照（RL 主干受控对比）。}A3C-R2N2 使用相同的 actor-critic 框架，但（a）编码器为 GRU+残差而非本文的 MLP，（b）状态与奖励为通用设计、不含 AIGC 成分。这是受控程度最高的对比。A3C-R2N2 取得 \slo${} = 0.214 \pm 0.07$、\etok${} = 2.93$~J/tok，本文为 $0.302 \pm 0.07$ 与 $2.61$~J/tok。差异在 \slo{} 维度 $p<0.01$、\etok{} 维度 $p<0.05$ 下显著。由于主干等价（均为带充足预训练的 actor-critic），这 $41\%$ 的 \slo{} 提升与 $11\%$ 的能效提升\textbf{可直接归因于 AIGC 感知的状态与奖励设计}。

\textbf{与 GNN 对照（状态编码受控对比）。}GNN 变体将 AIGC 物理编码为任务--服务器二部图中的\textbf{边特征}，是本文“MLP + 扁平 AIGC 特征”的替代方案。尽管架构更复杂，GNN 仅取得 \slo${} = 0.210 \pm 0.08$、\etok${} = 2.93$~J/tok，与 A3C-R2N2 基本持平，显著弱于本文调度器（$p<0.01$ \slo、$p<0.05$ \etok）。我们将该差距归因于\textbf{训练数据效率}：GNN 的关系推理能力需要比预训练预算（20 episode）更多的样本才能充分专门化。

\textbf{与 CS-UCB 对照（PerLLM 风格 bandit）。}CS-UCB 与本文共享逐请求粒度与能耗目标，其约束满足机制对应 PerLLM~\cite{yang2024perllm}。但它落在基线簇中（\slo{} 0.193、\etok{} 2.84）：\slo{} 显著低于本文调度器（$p{=}0.002$），能耗方向性更差（$p{=}0.066$）。原因是结构性的：平稳的逐 arm 均值无法跟踪\emph{状态相关}的代价---模型驻留与批处理槽位取决于集群不断演化的状态，上下文无关的 bandit 将其平均掉，而深度策略以其为条件。约束过滤不能替代基于状态的学习。

\textbf{与 NSGA-II 对照（多目标搜索）。}NSGA-II 是目前为止最强的基线，值得坦率分析。它在逐窗口分配上针对与本文相同的两个目标进行搜索，在典型工作点取得 \slo{} 0.289、\etok{} 2.79---本文调度器在两个均值上均更好（0.302、2.61），但在 $N{=}30$ 下两个差距均未单独显著（$p=0.64$、$p=0.12$）。三点观察补全图景。第一，在\textbf{中等负载甜区}中分离是真实的：$\lambda{=}1$ 处本文 \slo{} 达 0.447，NSGA-II 为 0.352（相对 $+27\%$，$p=0.009$）。第二，在我们运行 NSGA-II 的\textbf{全部八个配置中，本文每 token 能耗均更低}，优势从 $\lambda{=}0.5$ 处的持平增长到高负载下的 7--9\%。第三，两者在\textbf{调度时代价}上不对称：NSGA-II 在派发时刻重新运行种群搜索（总 wall-clock 为本文的 $5.4\times$；实测每派发一个任务约 70~ms 搜索），而本文策略经一次性预训练后以亚毫秒前向推理决策---在延迟敏感的服务场景中，派发时搜索恰恰会推高它所要保护的 TTFT。因此我们将 NSGA-II 视为最强的\emph{离线}竞争者；其在饱和工况的强势（\S\ref{sec:load}）使工况刻画（C4）更加清晰。

\subsection{拓扑敏感性}
\label{sec:topology}

在固定 $\lambda{=}2$ 请求/秒与 uniform 模型混合下，将边缘服务器数从 3 变到 7（始终配 1 台云端）。表~\ref{tab:topology} 报告 4 个最强基线与本文调度器的 \slo{} 与 \etok{}。

\begin{table}[t]
  \centering
  \caption{\slo{} 与 \etok{}（J/tok）随边缘数量变化。$\bigstar$ = 全部 11 种调度器中的列最优。}
  \label{tab:topology}
  \scriptsize
  \setlength{\tabcolsep}{2pt}
  \begin{tabular}{lrrrrrrr}
    \toprule
    配置（1 云 +）
      & PSO & SQ & A3C-R2N2 & GNN & CS-UCB & NSGA-II & \textbf{RL} \\
    \midrule
    3 边 --- \slo
      & 0.273 & 0.277 & 0.246 & 0.241 & 0.243 & 0.296
      & \textbf{0.314}$^\bigstar$ \\
    3 边 --- \etok
      & 2.19 & 2.32 & 2.28 & 2.24 & 2.24 & 2.16
      & \textbf{1.99}$^\bigstar$ \\
    5 边 --- \slo
      & 0.235 & 0.234 & 0.214 & 0.210 & 0.193 & 0.289
      & \textbf{0.302}$^\bigstar$ \\
    5 边 --- \etok
      & 2.81 & 2.96 & 2.93 & 2.93 & 2.84 & 2.79
      & \textbf{2.61}$^\bigstar$ \\
    7 边 --- \slo
      & 0.241 & 0.218 & 0.193 & 0.250 & 0.152 & 0.277
      & \textbf{0.282}$^\bigstar$ \\
    7 边 --- \etok
      & 3.25 & 3.34 & 3.35 & 3.35 & 3.35 & 3.19
      & \textbf{3.00}$^\bigstar$ \\
    \bottomrule
  \end{tabular}
\end{table}

在全部三种拓扑配置下，本文 RL 调度器在每项指标上均为最优。相对单目标基线，RL 的 \slo{} 相对领先 13--28\%、能耗领先 7--9\%。最强挑战者 NSGA-II 将 \slo{} 差距收窄到相对 2--6\%，但在每种拓扑下每 token 能耗都要多付 6--8\%---与我们在典型工作点观察到的不对称相同。

\textbf{能效优势比 \slo{} 更“配置稳定”。}相对能效提升在全部拓扑配置下都落在 [7\%, 9\%] 的窄带内，而 \slo{} 提升波动更大（13--28\%）。这表明\textbf{能效比 \slo{} 更干净地刻画了 AIGC 感知调度的“内在”收益}：\slo{} 对集群争用动态敏感，而能耗直接反映每任务效率。

\subsection{负载敏感性}
\label{sec:load}

固定 edge${=}5$，扫描 Poisson 到达率 $\lambda \in \{0.5, 1, 2, 4, 8\}$ 请求/秒，刻画调度器行为从欠利用到过饱和的变化。表~\ref{tab:slo-vs-lambda} 报告整个扫描范围内的 \slo{} 达成率与每 token 能耗。

\begin{table}[t]
  \centering
  \caption{负载扫描：\slo{} 达成率（上）与每 token 能耗（J/tok，下）随到达率变化。$\bigstar$ = 严格胜出；$\dagger$ = 与 PSO 持平。}
  \label{tab:slo-vs-lambda}
  \scriptsize
  \setlength{\tabcolsep}{2pt}
  \begin{tabular}{rrrrrrrl}
    \toprule
    $\lambda$ & PSO & SQ & A3C-R2N2 & GNN & CS-UCB & NSGA-II
      & \textbf{RL} \\
    \midrule
    \multicolumn{8}{@{}l}{\emph{\slo{} 达成率}} \\
    0.5 & 0.596 & 0.339 & 0.359 & 0.500 & 0.462 & 0.577
        & \textbf{0.596}$^\dagger$ \\
    1.0 & 0.297 & 0.265 & 0.227 & 0.237 & 0.214 & 0.352
        & \textbf{0.447}$^\bigstar$ \\
    2.0 & 0.235 & 0.234 & 0.214 & 0.210 & 0.193 & 0.289
        & \textbf{0.302}$^\bigstar$ \\
    4.0 & 0.253 & 0.205 & 0.191 & 0.207 & 0.215
        & \textbf{0.267} & 0.214 \\
    8.0 & 0.216 & 0.167 & 0.156 & 0.193 & 0.187
        & \textbf{0.219} & 0.199 \\
    \addlinespace
    \multicolumn{8}{@{}l}{\emph{每 token 能耗（J/tok）}} \\
    0.5 & 3.00 & 3.49 & 3.38 & 3.13 & 3.19 & 2.97
        & \textbf{2.96}$^\bigstar$ \\
    1.0 & 2.77 & 3.00 & 3.07 & 2.95 & 2.93 & 2.84
        & \textbf{2.64}$^\bigstar$ \\
    2.0 & 2.81 & 2.96 & 2.93 & 2.93 & 2.84 & 2.79
        & \textbf{2.61}$^\bigstar$ \\
    4.0 & 2.85 & 2.97 & 2.94 & 2.92 & 2.86 & 2.78
        & \textbf{2.58}$^\bigstar$ \\
    8.0 & 2.94 & 2.89 & 2.98 & 2.91 & 2.99 & 2.78
        & \textbf{2.52}$^\bigstar$ \\
    \bottomrule
  \end{tabular}
\end{table}

\textbf{关键观察。}（1）\emph{能效维度在所有负载下均为 RL 占优。}RL 在每个 $\lambda$ 都取得最低每 token 能耗（5/5），且相对次优基线的优势随负载单调增长---从 $\lambda{=}0.5$ 处的持平到 $\lambda{=}8$ 处的 $9\%$。（2）\emph{\slo{} 占优集中于中等负载工况}（$\lambda \in [0.5, 2]$）：在每一点 RL 均严格胜出或持平；$\lambda{=}1$ 处相对最强基线（NSGA-II，$p=0.009$）提升 $27\%$、相对 PSO 提升 $50\%$。（3）\emph{在极高负载下（$\lambda \ge 4$）}，集群逼近其物理吞吐上限，离线多目标搜索接管 \slo{} 轴：NSGA-II 在 $\lambda{=}4$ 与 $\lambda{=}8$ 处严格胜出，PSO 紧随其后---批量搜索比任何增量式策略更能利用同质化的队列状态---甚至 CS-UCB 也在 $\lambda{=}4$ 处达到 RL 的 \slo{} 水平（0.215 vs.\ 0.214）。\textbf{但 RL 的能效优势在该工况下持续存在且达到峰值}（表~\ref{tab:slo-vs-lambda} 下块）。

\textbf{AIGC 感知设计在什么工况下最有用？}我们刻画三个工况：\textbf{轻载}（$\lambda \le 1$，资源有余）---两轴均最优或持平；\textbf{甜区}（$1 \le \lambda \le 2$）---严格 Pareto 胜出，且 $\lambda{=}1$ 处的 \slo{} 领先即使相对 NSGA-II 也显著；\textbf{饱和}（$\lambda \ge 4$，吞吐上限）---仅保有能效优势，\slo{} 轴让位于离线多目标搜索。这一细致结论比扁平的“我们到处赢”更可辩护。最干净的表述是：\emph{AIGC 感知 RL 在每个负载下都取得最低每 token 能耗---且优势随负载增长---并在刻画多数生产推理部署的中等负载工况下取得严格的 \slo{}+能耗 Pareto 占优}。

\subsection{工作负载组合敏感性}
\label{sec:workload}

在 edge${=}5$、$\lambda{=}2$ 下测试两种模型混合：\textbf{uniform}（LLaMA-7B/13B/70B 各 1/3）与 \textbf{large-heavy}（70\% LLaMA-70B-INT8，偏向显存受限请求）。表~\ref{tab:workload} 报告结果。

\begin{table}[t]
  \centering
  \caption{不同工作负载混合下的 Pareto 表现。}
  \label{tab:workload}
  \small
  \begin{tabular}{llrr}
    \toprule
    Mix & Scheduler & \slo{} & \etok{} \\
    \midrule
    uniform & PSO            & 0.235 & 2.81 \\
    uniform & ShortestQueue  & 0.234 & 2.96 \\
    uniform & A3C-R2N2       & 0.214 & 2.93 \\
    uniform & GNN            & 0.210 & 2.93 \\
    uniform & CS-UCB         & 0.193 & 2.84 \\
    uniform & NSGA-II        & 0.289 & 2.79 \\
    uniform & \textbf{RL（本文）}
            & \textbf{0.302}$^\bigstar$ & \textbf{2.61}$^\bigstar$ \\
    \midrule
    large-heavy & PSO            & 0.093 & 3.84 \\
    large-heavy & ShortestQueue  & 0.105 & 3.84 \\
    large-heavy & A3C-R2N2       & 0.107 & 3.79 \\
    large-heavy & GNN            & 0.091 & 3.79 \\
    large-heavy & CS-UCB         & 0.068 & 3.84 \\
    large-heavy & NSGA-II        & \textbf{0.122} & 3.68 \\
    large-heavy & \textbf{RL（本文）}
                & 0.098 & \textbf{3.58}$^\bigstar$ \\
    \bottomrule
  \end{tabular}
\end{table}

在 \textbf{uniform 混合}下，本文 RL 为严格 Pareto 胜出者。在 \textbf{large-heavy 混合}下图景改变：所有调度器的 \slo{} 都坍缩到一个低带（0.07--0.12，uniform 下为 0.19--0.30），因为 LLaMA-70B-INT8（70~GB 权重）物理上\emph{只能装入云端服务器}。无论选择哪种调度器，云端队列都会饱和；在该带内，NSGA-II 的离线搜索取得最佳 \slo{}（0.122），呼应其在 \S\ref{sec:load} 中的饱和优势。\textbf{能效是唯一可广泛区分的轴}，RL 在该轴上仍然领先（3.58，对 NSGA-II 的 3.68 与其余基线的 3.79--3.84）。\emph{当工作负载在特定层级超出集群容量时，AIGC 感知调度退化为纯能效优化}（\S\ref{sec:discuss} 讨论）。

\subsection{组件消融}
\label{sec:ablation}

为隔离本文 RL 调度器的哪些组件贡献于其 Pareto 占优，我们系统消融 12 个组件：一次关闭一个、其余保持不变。所有消融使用 edge${=}5$、$\lambda{=}2$ 请求/秒、uniform 混合、$N{=}30$ 次运行。表~\ref{tab:ablation} 报告全部 12 个配置。

\textbf{消融分类。}四大类：\emph{物理层}（\texttt{no\_batching}）；\emph{PPO 内部}（\texttt{no\_action\_mask}、\texttt{no\_gae}、\texttt{no\_pretrain}、\texttt{no\_entropy}）；\emph{AIGC 奖励}（\texttt{no\_warm\_reward}、\texttt{no\_batch\_reward}、\texttt{no\_affinity\_reward}、\texttt{no\_cloud\_overuse}、\texttt{no\_all\_aigc\_rewards}）；\emph{AIGC 状态}（\texttt{no\_aigc\_state}、\texttt{no\_aigc\_full}）。

\begin{table*}[t]
  \centering
  \caption{组件消融。$\Delta$ 相对完整 RL（$N{=}30$、edge${=}5$、$\lambda{=}2$、uniform 混合）。$p$ 值来自对 \slo{} 相对完整 RL 的 Mann-Whitney U 检验。粗体 = $p<0.05$。}
  \label{tab:ablation}
  \small
  \begin{tabular}{llrrrrr}
    \toprule
    Category & Configuration & \slo{} & $\Delta$\slo\%
             & \etok{} & $\Delta$\etok\% & $p$(\slo) \\
    \midrule
    --- & \textbf{完整 RL（本文）}
        & \textbf{0.302} & ---
        & \textbf{2.61} & --- & --- \\
    \midrule
    \textbf{物理}
      & \textbf{$-$ 连续批处理}
      & \textbf{0.051} & \textbf{$-$83.2\%}
      & \textbf{5.75} & \textbf{$+$120.5\%}
      & \textbf{$<$0.001}$^{***}$ \\
    \textbf{PPO}
      & \textbf{$-$ 预训练（20$\to$0 ep）}
      & \textbf{0.186} & \textbf{$-$38.4\%}
      & \textbf{2.96} & \textbf{$+$13.5\%}
      & \textbf{$<$0.001}$^{***}$ \\
    PPO & $-$ 动作掩码
      & 0.274 & $-$9.3\% & 2.66 & $+$2.1\% & 0.39 \\
    PPO & $-$ GAE（改用 Monte-Carlo return）
      & 0.292 & $-$3.3\% & 2.61 & $-$0.1\% & 0.65 \\
    PPO & $-$ 熵正则
      & 0.305 & $+$0.9\% & 2.54 & $-$2.8\% & 0.92 \\
    AIGC 奖励 & $-$ Warm bonus
      & 0.300 & $-$0.7\% & 2.60 & $-$0.4\% & 0.87 \\
    AIGC 奖励 & $-$ Batch bonus
      & 0.311 & $+$3.1\% & 2.59 & $-$0.8\% & 0.87 \\
    AIGC 奖励 & $-$ Affinity bonus
      & 0.287 & $-$4.9\% & 2.64 & $+$1.3\% & 0.57 \\
    AIGC 奖励 & $-$ Cloud-overuse 惩罚
      & 0.309 & $+$2.2\% & 2.60 & $-$0.5\% & 0.96 \\
    AIGC 奖励（联合） & $-$ 所有 AIGC 奖励
      & 0.302 & $+$0.0\% & 2.62 & $+$0.5\% & 0.95 \\
    AIGC 状态 & $-$ AIGC 状态特征
      & 0.295 & $-$2.2\% & 2.55 & $-$2.1\% & 0.81 \\
    AIGC（联合） & $-$ 所有 AIGC 设计
      & 0.306 & $+$1.3\% & 2.52 & $-$3.3\% & 0.99 \\
    \bottomrule
  \end{tabular}
\end{table*}

\textbf{发现：两个主导组件与一个惊人的零结果。}

\textbf{发现 1：连续批处理是物理基础}（$p<0.001$、$\Delta$\slo${} = -83.2\%$、$\Delta$\etok${} = +120.5\%$）。关闭批处理使 \slo{} 坍缩到 5\%、每 token 能耗翻倍不止，印证了 Orca~\cite{yu2022orca} 与 Sarathi-Serve~\cite{agrawal2024sarathi} 将迭代级批处理视为基础性基础设施的强调。

\textbf{发现 2：充足预训练不可或缺}（$p<0.001$、$\Delta$\slo${} = -38.4\%$、$\Delta$\etok${} = +13.5\%$）。没有预训练，放置策略坍缩向云端服务器---与既有观察（Decima~\cite{mao2019decima}、RLScheduler~\cite{zhang2020rlscheduler}）一致：在随机调度环境中 warm-start 不可或缺。

\textbf{发现 3：单一 AIGC 组件无显著效应}（所有 $p > 0.39$）。关闭任何单个 AIGC 组件产生的变化都在噪声范围内（$\pm 5\%$）---与如下解释一致：在充足基础特征与预训练下，PPO 会\emph{隐式学到} AIGC 感知行为，显式整形提供的是冗余信号。

\textbf{发现 4（最惊人）：去掉\emph{全部} AIGC 设计同样没有影响}（\slo{} 的 $p=0.99$，$\Delta$\etok${} = -3.3\%$）。联合消融 \texttt{no\_aigc\_full} 的 \slo{} 表现与完整 RL 基本相同，能效甚至\emph{略好}。结合 \S\ref{sec:main-results} 中“本文调度器 Pareto 占优十个基线---但其 AIGC 专属设计组件却可被单独乃至联合消融”的发现，这描绘出一幅更微妙的图景：\textbf{“AIGC 感知”的贡献体现在训练流程 + 仿真器物理 + RL 主干选择之中，而非任何单一奖励或状态组件}（图~\ref{fig:ablation}）。

\begin{figure*}[t]
  \centering
  \includegraphics[width=0.88\textwidth]{fig7_ablation.pdf}
  \caption{12 个配置的组件消融。左：$\Delta$\slo\% vs.\ 完整 RL；右：$\Delta$\etok\%。红条 = $p<0.05$；灰条 = 噪声范围内；浅灰带 = $\pm 10\%$。只有连续批处理与预训练显示统计显著效应；所有单一 AIGC 奖励与状态组件均在噪声范围内。}
  \label{fig:ablation}
\end{figure*}

\textbf{调和：为何 RL 能 Pareto 占优基线却对消融不敏感？}\S\ref{sec:main-results}（RL 在均值上占优 10 个基线）与表~\ref{tab:ablation}（无单一 AIGC 组件显著）之间的表面矛盾可这样化解：每个基线与本文调度器同时在\emph{多个}根本维度上不同（无学习；固定代价函数且不含 AIGC 物理；不同的编码器或训练配方）---因而落后 7--28\%。而当我们在本文调度器内部消融任何\emph{单个} AIGC 组件时，PPO 会从剩余特征中提取等价信号加以补偿。这告诉我们：\textbf{AIGC 感知调度是一项整体性的系统贡献，而非某个巧妙的奖励或状态组件}。

\subsection{奖励权重敏感性}
\label{sec:weight-sens}

式~\ref{eq:reward} 的奖励把一个多目标问题坍缩为标量加权和---我们的结论是否只是某个特定权重向量的产物？我们在覆盖设计空间的四种替代配置下从零重训---\emph{uniform}（全部权重 $1/7$）、\emph{time-heavy}（时间项占一半以上权重）、\emph{AIGC-heavy}（warm{+}batch{+}affinity 提升到 0.70）与 \emph{base-heavy}（AIGC 项几乎移除）---每种配置在典型工作点运行 $N{=}30$ 次（表~\ref{tab:weight-sens}）。

\begin{table}[t]
  \centering
  \caption{奖励权重敏感性（edge${=}5$、$\lambda{=}2$~请求/秒、每配置 $N{=}30$）。权重按式~\ref{eq:reward} 的 time/balance/match/warm/batch/affinity/cloud 顺序列出。$p$：对典型配置的双侧 Mann-Whitney U，\slo{} / \etok。每个配置均 Pareto 占优最强基线（PSO：\slo{}~0.235、\etok{}~2.81）。}
  \label{tab:weight-sens}
  \scriptsize
  \setlength{\tabcolsep}{2.5pt}
  \begin{tabular}{@{}llccc@{}}
    \toprule
    Config & Weights & \slo{} & \etok{} & $p$ \\
    \midrule
    canonical & .25/.10/.10/.15/.20/.10/.10
      & $0.302 {\pm} 0.13$ & $2.61 {\pm} 0.67$ & --- \\
    uniform & 各 $1/7$
      & $0.287 {\pm} 0.13$ & $2.61 {\pm} 0.68$ & .72 / .98 \\
    time-heavy & .55/.05/.05/.10/.10/.05/.10
      & $0.304 {\pm} 0.14$ & $2.60 {\pm} 0.67$ & .97 / .98 \\
    AIGC-heavy & .10/.05/.05/.25/.30/.15/.10
      & $0.300 {\pm} 0.12$ & $2.63 {\pm} 0.67$ & .86 / .96 \\
    base-heavy & .40/.20/.20/.05/.05/.05/.05
      & $0.305 {\pm} 0.14$ & $2.56 {\pm} 0.68$ & .89 / .74 \\
    \bottomrule
  \end{tabular}
\end{table}

两点观察。第一，\textbf{核心结论对权重选择稳健}：每个配置---包括极端配置---都 Pareto 占优最强基线（\slo{} 0.287--0.305 vs.\ PSO 的 0.235；\etok{} 2.56--2.63 vs.\ 2.81）。第二，\textbf{没有任何配置与典型权重存在显著差异}（所有 $p \ge 0.72$）：即便把 AIGC 组的权重质量从 0.15（base-heavy）移到 0.70（AIGC-heavy），两个目标在统计上均保持不变。这把消融的零结果从\emph{移除}组件扩展到\emph{重加权}组件，并说明我们的标量化是良性的：在权重空间的一个宽阔区域内，学到的策略落在同一个 Pareto 点上，因此从典型配置得出的结论具有普适性。
